% ---------------------------------------------------------------- 3.3 aside: why the sigmoid
\begin{frame}{Why a Sigmoid? Why Not Just Subtract the Scores?}
\textbf{Because we need a probability, and the gap is not one.}
\vspace{0.6em}
\begin{itemize}\itemsep0.45em\small
    \item Each label is a \emph{binary event}: the human clicked $y^+$. Fitting by
          \textbf{maximum likelihood} means choosing $\phi$ to make the clicks we
          actually observed as likely as possible under the model, i.e.\ maximizing
          \[
              \prod_{\text{comparisons}} P_\phi(\text{the choice the human made})
          \]
    \item \empr{Every factor in that product has to be a probability.} If the model
          emits an unbounded number instead, the product is not a likelihood of
          anything: it has no maximum (push the number up forever), and its $\log$ is
          undefined once it goes negative. There is nothing to fit.
    \item So the model must output $P\in[0,1]$ per comparison, while the score gap
          $r(y^+)-r(y^-)$ is a free real number. $\sigma$ is the monotone map
          $\mathbb{R}\to(0,1)$ with the right endpoints: gap $0 \to 0.5$
          (a coin flip, the model is indifferent), gap $\to+\infty \to 1$,
          gap $\to-\infty \to 0$.
\end{itemize}
\end{frame}

